\titleformat{\section}
  {\normalfont\large\bfseries}
  {\appendixname~\thesection.}{0.5em}{}

\section{Proof of Theorem~\ref{thm:single-cell-representation}}
\label{sec:appendixA}
In this appendix, the superscripts ($+$) and ($-$) no longer label the right
and left half-guides, respectively. Instead, they denote the exterior and
interior traces on \(\partial\Omega\), respectively. Let \(\nu\) be the unit
normal on \(\partial\Omega\) pointing out of \(\Omega\). We define
\begin{align*}
  u^\pm(\mathbf{x})
  &:= \lim_{h \to 0^+}u\bigl(\mathbf{x}\pm h\nu(\mathbf{x})\bigr),\\
  u_\nu^\pm(\mathbf{x})
  &:= \lim_{h \to 0^+}
  \nabla u\bigl(\mathbf{x}\pm h\nu(\mathbf{x})\bigr)\cdot\nu(\mathbf{x}),
  \qquad \mathbf{x}\in\partial\Omega.
\end{align*}

We first recall the free-space Green representation formulas used below; see
\cite[Section~3.2]{ColtonKress1983}. Let
\(\kappa>0\), and suppose that \(u\) is a sufficiently smooth solution of
\[
  (\Delta+\kappa^2)u=0 \qquad \text{in }\Omega.
\]
Then its interior Cauchy data satisfy
\begin{align}
  -\mathcal{S}^{(\kappa)}[u_\nu^-](\mathbf{x})
  +\mathcal{D}^{(\kappa)}[u^-](\mathbf{x})
  =
  \begin{cases}
    -u(\mathbf{x}), & \mathbf{x}\in\Omega,\\
    0, & \mathbf{x}\in\mathbb{R}^2\setminus\overline{\Omega}.
  \end{cases}
  \label{eq:green-representation-interior}
\end{align}
The second line is the corresponding extinction formula: the layer potential
generated by the interior Cauchy data vanishes on the other side of the
interface.

For the exterior formula, suppose instead that \(u\) solves
\[
  (\Delta+\kappa^2)u=0
  \qquad \text{in }\mathbb{R}^2\setminus\overline{\Omega}
\]
and satisfies the two-dimensional Sommerfeld radiation condition
\begin{align}
  \frac{\partial u}{\partial r}-\ii\kappa u
  =o\bigl(r^{-1/2}\bigr),
  \qquad r=|\mathbf{x}|\to\infty,
  \label{eq:sommerfeld-radiation-appendix}
\end{align}
uniformly with respect to the direction \(\mathbf{x}/|\mathbf{x}|\). Its
exterior Cauchy data then give
\begin{align}
  -\mathcal{S}^{(\kappa)}[u_\nu^+](\mathbf{x})
  +\mathcal{D}^{(\kappa)}[u^+](\mathbf{x})
  =
  \begin{cases}
    0, & \mathbf{x}\in\Omega,\\
    u(\mathbf{x}), & \mathbf{x}\in\mathbb{R}^2\setminus\overline{\Omega}.
  \end{cases}
  \label{eq:green-representation-exterior}
\end{align}
Thus the same combination \(-\mathcal{S}^{(\kappa)}u_\nu+
\mathcal{D}^{(\kappa)}u\) reproduces the Helmholtz field on the side from
which the Cauchy data are taken and vanishes on the opposite side, with the
sign determined by whether the data are interior or exterior traces.

We will need the following quasiperiodic analogues.

\begin{lemma}
\label{lem:qp-interior-green-representation}
Let $u$ satisfy $(\Delta+\kappa^2)u=0$ in $\Omega$, then for $(\kappa, \beta)$ not a Wood anomaly, we have\begin{align}
  -\mathcal{S}^{(\kappa)}_\mathrm{QP}[u_\nu^-](\mathbf{x})
  +\mathcal{D}^{(\kappa)}_\mathrm{QP}[u^-](\mathbf{x})
  =
  \begin{cases}
    -u(\mathbf{x}), & \mathbf{x}\in\Omega,\\
    0, & \mathbf{x}\in \mathcal{C}\setminus\overline{\Omega}.
  \end{cases}
\end{align}
\end{lemma}

\begin{proof}
Following the lattice-image argument of
\cite[Appendix~A, Lemma~9]{BarnettGreengard2010}, let
\(\mathbf e_y=(0,1)\), and for each \(j\in\mathbb Z\) define the translated
inclusion and field
\[
  \Omega_j:=\Omega+jd\mathbf e_y,
  \qquad
  u_j(\mathbf z):=u(\mathbf z-jd\mathbf e_y),
  \quad \mathbf z\in\Omega_j.
\]
Then \((\Delta+\kappa^2)u_j=0\) in \(\Omega_j\). By translation invariance of
the free-space Green's function, the lattice-sum definition of the
quasiperiodic potentials gives
\begin{align*}
  &-\mathcal S_\mathrm{QP}^{(\kappa)}[u_\nu^-](\mathbf x)
  +\mathcal D_\mathrm{QP}^{(\kappa)}[u^-](\mathbf x)\\
  &\qquad=
  \lim_{N\to\infty}\sum_{j=-N}^{N}e^{\ii j\beta d}
  \left(
    -\mathcal S_{\partial\Omega_j}^{(\kappa)}[(u_j)_\nu^-](\mathbf x)
    +\mathcal D_{\partial\Omega_j}^{(\kappa)}[u_j^-](\mathbf x)
  \right),
\end{align*}
where the boundary written as a subscript indicates the curve over which the
free-space potential is integrated.

Since the inclusion lies strictly inside the cell, every
\(\mathbf x\in\mathcal C\) lies outside \(\overline{\Omega_j}\) when
\(j\neq0\). The extinction part of
\eqref{eq:green-representation-interior}, applied to \(u_j\) in \(\Omega_j\),
therefore shows that every \(j\neq0\) term in the lattice sum is zero. For
\(j=0\), the same formula gives \(-u(\mathbf x)\) when
\(\mathbf x\in\Omega\) and zero when
\(\mathbf x\in\mathcal C\setminus\overline\Omega\). Thus the only
contribution surviving the lattice-sum limit is the \(j=0\) image, which
proves the stated formula.
\end{proof}


\begin{lemma}
\label{lem:qp-exterior-green-representation}
Let $u$ satisfy $(\Delta+\kappa^2)u=0$ in $\mathcal{C}\setminus\overline{\Omega}$ and the quasiperiodic condition \begin{align*}
  \left. u \right|_{y = d} = e^{\ii \beta d} \left. u \right|_{y = 0}, \quad
\left. \partial_y u \right|_{y = d} = e^{\ii \beta d} \left. \partial_y u \right|_{y = 0}.
\end{align*}
Assume that $(\kappa, \beta)$ is not a Wood anomaly and that \(u\) satisfies the Rayleigh radiation condition, then
\begin{align}
  -\mathcal{S}^{(\kappa)}_\mathrm{QP}[u_\nu^+](\mathbf{x})
  +\mathcal{D}^{(\kappa)}_\mathrm{QP}[u^+](\mathbf{x})
  =
  \begin{cases}
    0, & \mathbf{x}\in\Omega,\\
    u(\mathbf{x}), & \mathbf{x}\in \mathcal{C}\setminus\overline{\Omega}.
  \end{cases} \label{eq:qp-exterior-green-representation}
\end{align}
\end{lemma}

\begin{proof}
We follow the standard method of proof \cite[Theorem~3.3]{ColtonKress1983}. Apply Green's 2nd identity to the domain $\mathcal{C} \setminus \overline{\Omega}$ if 
$\mathbf{x} \in \Omega$, or the domain $\{\mathbf{y} \in \mathcal{C}\setminus \overline{\Omega} : |\mathbf{x} - \mathbf{y}| > \varepsilon\}$ if $\mathbf{x} \in \mathcal{C} \setminus \overline{\Omega}$. In the latter case the limit $\varepsilon \to 0$ is taken, 
and \eqref{eq:lattice-sum-qpgreen} shows that only $j = 0$ term contributes to the limit
of the integral over the circle of radius $\varepsilon$ centered at $\mathbf{x}$. The 
integral over the boundary $\partial \Omega$ gives the layer potential combination.
We will focus on the following integral:\begin{align}
  I = \int_{\partial \mathcal{C}} \frac{\partial G^{(\kappa)}_\mathrm{QP}(\mathbf{x},\mathbf{y})}{\partial \nu(\mathbf{y})} u(\mathbf{y}) - G^{(\kappa)}_\mathrm{QP}(\mathbf{x},\mathbf{y}) u_\nu(\mathbf{y}) \, \dd s(\mathbf{y}).
\end{align}
The integrals over the top and bottom sides of the cell cancel because $G^{(\kappa)}_\mathrm{QP}(\mathbf{x}, \cdot)$ is anti-quasiperiodic, i.e. quasiperiodic with phase 
$e^{-\ii \beta d}$. Suppose $\mathbf{y} = (s,t)$, then we have \begin{align*}
  I_\mathrm{top} &= \int_{X_\mathrm{L}}^{X_\mathrm{R}} \left[e^{\ii \beta d} u(s,0) e^{-\ii \beta d}\partial_t G^{(\kappa)}_\mathrm{QP}(\mathbf{x},(s,0)) - e^{-\ii \beta d} G^{(\kappa)}_\mathrm{QP}(\mathbf{x},(s,0)) e^{\ii \beta d} \partial_t u(s,0) \right] \, \dd s,\\
  I_\mathrm{bottom} &= -\int_{X_\mathrm{L}}^{X_\mathrm{R}} \left[u(s,0) \partial_t G^{(\kappa)}_\mathrm{QP}(\mathbf{x},(s,0)) - G^{(\kappa)}_\mathrm{QP}(\mathbf{x},(s,0))\partial_t u(s,0)\right] \, \dd s,
\end{align*}
which cancel each other. 

The integrals over the left and right sides need to be evaluated carefully. By\eqref{eq:right-coefficient-extraction}--\eqref{eq:right-outgoing-rayleigh-condition},
the Rayleigh radiation condition eliminates the
right-going Rayleigh component at the left boundary, and the left-going Rayleigh component at the right boundary, which implies that $u$ has the form 
\begin{align*}
  u(s,t) &= \sum_{m\in\mathbb Z} a_m^\mathrm{R} e^{\ii\gamma_m (s - X_\mathrm{R})}\psi_m(t), \quad s \geq X_\mathrm{R},\\
  u(s,t) &= \sum_{m\in\mathbb Z} b_m^\mathrm{L} e^{-\ii\gamma_m (s - X_\mathrm{L})}\psi_m(t), \quad s \leq X_\mathrm{L},
\end{align*}
where $a, b$ are given by \eqref{eq:right-coefficient-extraction}--\eqref{eq:left-coefficient-extraction}. Define the anti-quasiperiodic dual modes 
\begin{align*}
  \widetilde{\psi}_m(t) = \frac{1}{\sqrt{d}} e^{-\ii \beta_m t},\quad \int_0^d \psi_m(t) \widetilde{\psi}_n(t) \, \dd t = \delta_{mn}.
\end{align*}
The spectral Green's function formula \eqref{eq:qp-green-spectral}, now viewed as a 
function of $\mathbf{y} = (s,t)$, gives: \begin{align*}
  G^{(\kappa)}_\mathrm{QP}(\mathbf{x},(s,t)) &= \sum_{m\in\mathbb Z} 
  c_m^\mathrm{L}(\mathbf{x}) e^{-\ii \gamma_m (s - X_\mathrm{L})} \widetilde\psi_m(t), \quad s < x_1,\\
  G^{(\kappa)}_\mathrm{QP}(\mathbf{x},(s,t)) &= \sum_{m\in\mathbb Z} 
  c_m^\mathrm{R}(\mathbf{x}) e^{\ii \gamma_m (s - X_\mathrm{R})} \widetilde\psi_m(t), \quad s > x_1,
\end{align*}
where, explicitly, \begin{align*}
  c_m^\mathrm{L}(\mathbf{x}) = \frac{\ii}{2\sqrt{d}\gamma_m} e^{\ii \beta_m x_2} 
  e^{\ii \gamma_m (x_1 - X_\mathrm{L})},\quad 
  c_m^\mathrm{R}(\mathbf{x}) = \frac{\ii}{2\sqrt{d}\gamma_m} e^{\ii \beta_m x_2}
  e^{\ii \gamma_m (X_\mathrm{R} - x_1)}.
\end{align*}
On the right side $s = X_\mathrm{R}$, we have \begin{align*}
  I_\mathrm{right} &= \int_0^d u(X_\mathrm{R},t) \partial_s G^{(\kappa)}_\mathrm{QP}(\mathbf{x},(X_\mathrm{R},t)) - G^{(\kappa)}_\mathrm{QP}(\mathbf{x},(X_\mathrm{R},t))\partial_s u(X_\mathrm{R},t)  \, \dd t\\
  &= \sum_{m,n} a_m^\mathrm{R} c_n^\mathrm{R}(\mathbf{x}) \ii(\gamma_n - \gamma_m) \int_0^d \psi_m(t) \widetilde{\psi}_n(t) \, \dd t = 0.
\end{align*}
An identical calculation gives $I_\mathrm{left} = 0$. Thus the integral over the boundary of the cell vanishes, the representation formula \eqref{eq:qp-exterior-green-representation} holds, and the proof is complete.
\end{proof}

Now we are ready to prove Theorem~\ref{thm:single-cell-representation}.
\begin{proof}[Proof of Theorem~\ref{thm:single-cell-representation}]
We prove the existence and uniqueness seperately.

\emph{1. Existence.}
In a homogeneous neighborhood of $\Gamma_\mathrm{L}$ not intersecting $\overline{\Omega}$,
field $u$ has unique Rayleigh expansion \begin{align*}
  u(x, y) = \sum_{m \in \mathbb{Z}} \left[ a_m^\mathrm{L}e^{\ii\gamma_m (x - X_\mathrm{L})} + b_m^\mathrm{L}e^{-\ii\gamma_m (x - X_\mathrm{L})} \right] \psi_m(y), 
\end{align*}
the Rayleigh coefficients $a_m^\mathrm{L}$ and $b_m^\mathrm{L}$ are given by \eqref{eq:right-coefficient-extraction}--\eqref{eq:left-coefficient-extraction} since $(k, \beta)$ is not a Wood anomaly.
In the neighborhood of $\Gamma_\mathrm{R}$ we can extract the Rayleigh coefficients $a_m^\mathrm{R}$ and $b_m^\mathrm{R}$ in a similar way. We define the homogeneous 
background component $u_h$ as following:
\begin{align}
  u_h(x, y) = \sum_{m \in \mathbb{Z}} \left[ a_m^\mathrm{L}e^{\ii\gamma_m (x - X_\mathrm{L})} + b_m^\mathrm{R}e^{-\ii\gamma_m (x - X_\mathrm{R})} \right] \psi_m(y),
\end{align}
and define \begin{align}
  w = \begin{cases}
    u - u_h, & \mathbf{x} \in \mathcal{C}\setminus\overline{\Omega},\\
    u, & \mathbf{x} \in \Omega,
  \end{cases}
\end{align}
satisfying the Rayleigh radiation condition \eqref{eq:left-outgoing-rayleigh-condition}--\eqref{eq:right-outgoing-rayleigh-condition}. 


We introduce the (swapped-wavenumber) transmission problem: \begin{align}
  \begin{cases}
  (\Delta + k^2) v = 0& \quad \text{in } \Omega,\\
  (\Delta + n^2k^2) v = 0& \quad \text{in } \mathbb{R}^2 \setminus \overline{\Omega},\\
  \frac{\partial v}{\partial r} - \ii n k v = o(r^{-1/2}),& \quad r = |\mathbf{x}| \to \infty,\\
  v^+ - v^- = f,&\\
  v_\nu^+ - v_\nu^- = g,&
  \end{cases}
  \label{eq:swapped-transmission-problem}
\end{align}
where $f$ and $g$ are given continuous functions on $\partial \Omega$, in our construction 
we take
\begin{align*}
  f = -w^+ - w^-; \qquad g = -w_\nu^+ - w_\nu^-,
\end{align*}
then the transmission problem has a unique solution $v$ by \cite[Theorem~3.41]{ColtonKress1983}.
Define the densities\begin{align*}
  \tau = w^+ - v^- = -w^- - v^+;\qquad \sigma = -w^+_\nu + v_\nu^- = w^-_\nu + v_\nu^+,
\end{align*}
It remains to verify that these densities produce the prescribed fields $w$. 
For $\mathbf{x} \in \Omega$, the Green's representation formulae \eqref{eq:green-representation-interior} and \eqref{eq:green-representation-exterior} give
\begin{align*}
  &\mathcal{S}^{(nk)}[\sigma](\mathbf{x}) + \mathcal{D}^{(nk)}[\tau](\mathbf{x})\\
  ={}& -\left[ -\mathcal{S}^{(nk)}[w_\nu^-](\mathbf{x}) + \mathcal{D}^{(nk)}[w^-](\mathbf{x}) \right] - \left[ -\mathcal{S}^{(nk)}[v_\nu^+](\mathbf{x}) + \mathcal{D}^{(nk)}[v^+](\mathbf{x}) \right]\\
  ={}& -[-w(\mathbf{x})] - 0 = w(\mathbf{x}), \qquad \mathbf{x} \in \Omega.
\end{align*}
For $\mathbf{x} \in \mathcal{C} \setminus \overline{\Omega}$, Lemma~\ref{lem:qp-interior-green-representation} and Lemma~\ref{lem:qp-exterior-green-representation} give
\begin{align*}
  &\mathcal{S}^{(k)}_\mathrm{QP}[\sigma](\mathbf{x}) + \mathcal{D}^{(k)}_\mathrm{QP}[\tau](\mathbf{x})\\
  ={}& \left[ -\mathcal{S}^{(k)}_\mathrm{QP}[w_\nu^+](\mathbf{x}) + \mathcal{D}^{(k)}_\mathrm{QP}[w^+](\mathbf{x}) \right] - \left[ -\mathcal{S}^{(k)}_\mathrm{QP}[v_\nu^-](\mathbf{x}) + \mathcal{D}^{(nk)}_\mathrm{QP}[v^-](\mathbf{x}) \right]\\
  ={}& w(\mathbf{x}) - 0 = w(\mathbf{x}), \qquad \mathbf{x} \in \mathcal{C} \setminus \overline{\Omega}.
\end{align*}
This proves existence.

\emph{2. Uniqueness.}
Suppose that
\[
  u=u_h+w=\widehat u_h+\widehat w
  \qquad\text{in }\mathcal C\setminus\overline\Omega
\]
are two decompositions of the stated form. Then
\[
  \widehat u_h-u_h=w-\widehat w.
\]
The right-hand side is outgoing at both ports. The left-hand side is a
homogeneous background field of the form \eqref{eq:homogeneous-field}, whose
only incoming coefficients are
$\widehat\xi_m^\rightarrow-\xi_m^\rightarrow$ at the left port and
$\widehat\xi_m^\leftarrow-\xi_m^\leftarrow$ at the right port. Since it is
also outgoing, all these coefficients vanish. Hence
$\widehat u_h=u_h$ and $\widehat w=w$. The interface densities are then
unique by the same complementary-field uniqueness argument used in their
construction above.

To prove the uniqueness of the densities $\sigma, \tau$, suppose that $w = 0$ in 
$\mathcal{C} \setminus \partial \Omega$, and define the complementary field over 
the whole plane minus $\partial \Omega$
\begin{align}
  v(\mathbf{x}) = \begin{cases}
    \mathcal{S}^{(k)}_\mathrm{QP}[\sigma](\mathbf{x}) + \mathcal{D}^{(k)}_\mathrm{QP}[\tau](\mathbf{x}), & \mathbf{x} \in \Omega,\\
    -\mathcal{S}^{(nk)}[\sigma](\mathbf{x}) - \mathcal{D}^{(nk)}[\tau](\mathbf{x}), & \mathbf{x} \in \mathbb{R}^2 \setminus \overline{\Omega}.
  \end{cases}
\end{align}
Then $w^- = w^-_\nu = 0$ and by the jump relations of 
$\mathcal{S}^{(nk)}[\sigma]+ \mathcal{D}^{(nk)}[\tau]$ we have $v^+ = -\tau$ and $v_\nu^+ = \sigma$. Similarly, since $w^+ = w^+_\nu = 0$ by the jump relations of $\mathcal{S}^{(k)}_\mathrm{QP}[\sigma]+ \mathcal{D}^{(k)}_\mathrm{QP}[\tau]$ we have $v^- = -\tau$ and $v_\nu^- = \sigma$.
It's easy to check that $v$ solves the transmission problem 
\eqref{eq:swapped-transmission-problem} with $f = g = 0$. By the uniqueness of the transmission problem, we have $v \equiv 0$ in $\mathbb{R}^2 \setminus \partial \Omega$, 
from which the jump relations back to the $w$ imply $\sigma = \tau = 0$. This proves the uniqueness of the densities.
\end{proof}
